\chapter{критический аудит-аудит физического динамического замыкания}
\label{ch:v8-dynamic-physical-closure-redteam-gate}

\section{Проверяемая граница}

К настоящему этапу построены типизированный конечномерный GKSL-генератор,
примитивная полугруппа, безразмерное шумовое время и общий цепной оператор
боровских частот. Эти результаты образуют содержательно замкнутый класс
операторных процессов. Они, однако, ещё не доказывают, что один из процессов
является единственной физической динамикой.

Для такого усиления потребуем одновременно:
\begin{enumerate}
 \item единственную относительную метрику скоростей;
 \item физическую калибровку энергии и времени;
 \item микроскопическую связь с состоянием среды или её спектральной
 плотностью;
 \item независимо заданное полное корреляционное ядро;
 \item вывод шумовой метрики из физической симметрии;
 \item заранее зарегистрированную наблюдаемую величину.
\end{enumerate}

Это усиленная динамическая версия правила старого глобального
проверки фальсифицируемости: внутренне точная формула не повышается до физического
предсказания без независимого отображения на наблюдаемое.

\section{Точный аудит ориентации и масштаба}

Цепной оператор даёт две KMS-ветви
\begin{equation}
 r_+=e^{-2},\qquad r_-=e^2,\qquad r_+r_-=1.
 \label{eq:v8-redteam-reciprocal-branches}
\end{equation}
Индексный гейт выбирает $r_+$ как ветвь, совместимую с уже фиксированным
знаком Ходж-функционала. Это строгий внутренний результат. Но он не
превращает знак градуировки в независимо измеренную термодинамическую стрелу
и не задаёт размерную величину $\beta\Delta$.

Аналогично для любого $\kappa>0$
\begin{equation}
 \ker(\kappa\mathcal L)=\ker\mathcal L.
 \label{eq:v8-redteam-rate-rescaling}
\end{equation}
Следовательно, неподвижная алгебра, примитивность и направленное отношение
скоростей не определяют перевод безразмерного параметра полугруппы в
физическое время.

\section{Сверка цепи происхождения}

Машинный аудит читает одиннадцать уже существующих сертификатов. Он
подтверждает положительное математическое ядро: полную положительность,
скалярную неподвижную алгебру, безразмерную полугруппу, общий боровский
разрыв и совместимую с родителем ориентацию. Одновременно он воспроизводит
восемь незакрытых физических требований:
\begin{enumerate}
 \item относительная скоростная метрика не единственна;
 \item абсолютный масштаб времени не выведен;
 \item физическая энергия или температура не откалибрована;
 \item автономная подача среды не построена;
 \item микроскопическая система--среда связь отсутствует;
 \item независимое физическое корреляционное ядро отсутствует;
 \item шумовая метрика не вынуждена симметрией;
 \item пререгистрированная наблюдаемая не предъявлена.
\end{enumerate}

Стандартный слабосвязанный предел строит марковский генератор из данных
системы, среды и их взаимодействия; KMS подробный баланс связывает уже
заданные направленные скорости, но не создаёт отсутствующую спектральную
плотность. Повторно-взаимодействующий предел также требует явного взаимодействия
и правила масштабирования. Поэтому обнаруженный дефицит является не
технической особенностью аудита, а правильной границей физической
интерпретации.

\begin{theorem}[Операторное замыкание без физического замыкания]
Текущая архитектура Тома VIII задаёт строгий конечномерный класс
примитивных квантовых марковских процессов с типизированными скачками,
безразмерным временем и внутренне согласованной ориентацией. Она не задаёт
единственный физический генератор, поскольку не выводит относительную
скоростную метрику, размерный масштаб, микроскопическую среду и независимое
корреляционное ядро.
\end{theorem}

\begin{nogocriterion}[Запрет погони за внутренними селекторами]
Нельзя повышать очередной канонический представитель, индексный знак,
следовый Казимир или Ходж-вес до физического закона только потому, что он
уменьшает пространство свободных параметров. Новое физическое замыкание
разрешено лишь после предъявления микроскопической система--среда модели либо
пререгистрированной наблюдаемой, различающей оставшиеся процессы.
\end{nogocriterion}

\begin{protocol}
\begin{enumerate}
 \item типизированный конечномерный генератор: \textbf{получен};
 \item полная положительность и примитивность: \textbf{получены};
 \item безразмерное время: \textbf{получено};
 \item внутренняя Ходж-совместимая ориентация: \textbf{получена};
 \item физическая стрела времени: \textbf{не выведена независимо};
 \item единственная относительная скоростная метрика: \textbf{не получена};
 \item абсолютная физическая скорость: \textbf{не получена};
 \item микроскопическая среда или спектр среды: \textbf{не предъявлены};
 \item независимое физическое ядро: \textbf{не предъявлено};
 \item математическая программа: \textbf{не находится в тупике};
 \item уникальное физическое динамическое замыкание: \textbf{запрет};
 \item допустимое продолжение: \textbf{LCF-аудит без повышения физического статуса};
 \item следующий физический гейт: \textbf{система--среда спектральная плотность
 или пререгистрированная наблюдаемая}.
\end{enumerate}
\end{protocol}

Машинный сертификат сохранён в
\path{s2t_v8_dynamic_physical_closure_redteam_gate_results.json}.

\begin{thebibliography}{9}
\bibitem{v8-redteam-davies}
E.~B.~Davies,
\emph{Markovian master equations},
Communications in Mathematical Physics 39 (1974), 91--110.
\bibitem{v8-redteam-fagnola}
F.~Fagnola, V.~Umanit\`a,
\emph{Generators of detailed balance quantum Markov semigroups},
arXiv:0707.2147.
\bibitem{v8-redteam-attal}
S.~Attal, Y.~Pautrat,
\emph{From repeated to continuous quantum interactions},
Annales Henri Poincar\'e 7 (2006), 59--104.
\bibitem{v8-redteam-connes}
A.~Connes, C.~Rovelli,
\emph{Von Neumann algebra automorphisms and time--thermodynamics relation
in generally covariant quantum theories},
Classical and Quantum Gravity 11 (1994), 2899--2918.
\end{thebibliography}